% Copyright (c) 2026 Benjamin Stanley Frohman. All rights reserved.
% Released under Apache 2.0 license as described in the file LICENSE.
% Authors: Benjamin Stanley Frohman
\documentclass[11pt]{amsart}
\usepackage{amsmath,amssymb,amsthm,booktabs,hyperref}
\usepackage[margin=1.15in]{geometry}

\title{The Type of a Rational Hodge Counterexample}
\author{Benjamin Stanley Frohman}
\address{Independent Researcher \\ ORCID: 0009-0006-7068-3718}
\email{frohmanbenjamin@gmail.com}
\date{September 24, 2026}

\begin{document}

\begin{abstract}
The Clay Millennium statement of the Hodge conjecture is the \emph{rational} sentence
\[
  \forall D\,\forall\gamma\,\exists z,\quad \mathrm{cl}(z)=\gamma,
\]
with coefficients in $\mathbb{Q}$. This note records the two Lean~4 types of its negation, proves they are classically equivalent, and records that both types are uninhabited in the accompanying formalization (\texttt{BenFrohman/HodgeDisproof}). It does not produce a fourfold $D_{\mathrm{bad}}$ or a class $\gamma_{\mathrm{bad}}$. The integral Hodge conjecture is a different sentence and is already false.
\end{abstract}

\maketitle

\section{The Sentence}

Let $D$ range over smooth complex projective fourfolds and let $\gamma$ range over rational Hodge classes of type $(2,2)$:
\[
  \gamma\in H^4(D,\mathbb{Q})\cap H^{2,2}(D).
\]
Let $z$ range over rational algebraic cycles of codimension two and let $\mathrm{cl}$ be the cycle class map. The easy arrow---every such $\mathrm{cl}(z)$ is Hodge---is a theorem. The Clay problem is the converse, expressed by the universal sentence above.

Hodge's original integral sentence, with $\mathbb{Z}$ in place of $\mathbb{Q}$, is false: Atiyah--Hirzebruch (1961) produced torsion classes that are not algebraic; Koll\'ar (1990) produced infinite-order integral classes that are not algebraic. Tensoring with $\mathbb{Q}$ kills torsion and identifies classes that differ by an integer multiple. Those terms do not inhabit the rational negation.

\section{Two Types of Negation}

In dependent type theory the constructive negation of a universal sentence is a function type, whereas a counterexample witness is a dependent pair (sigma/existential triple).

\begin{center}
\renewcommand{\arraystretch}{1.25}
\begin{tabular}{@{}lll@{}}
\toprule
\textbf{Lean Name} & \textbf{Type} & \textbf{Inhabitant} \\
\midrule
\texttt{RationalHodgeNegation} & $(\forall D\,\forall\gamma\,\exists z,\,\mathrm{cl}\,z=\gamma)\to\mathrm{False}$ & a function \\
\texttt{RationalCounterexample} & $\exists D\,\exists\gamma,\,\forall z,\,\mathrm{cl}\,z=\gamma\to\mathrm{False}$ & a triple $\langle D,\gamma,p\rangle$ \\
\bottomrule
\end{tabular}
\end{center}

Lean treats these as distinct types; the kernel will not accept a triple where it expects a function.

\paragraph{\textbf{Constructive map (\texttt{counterexample\_to\_negation}).}}
From a triple $\langle D_{\mathrm{bad}},\gamma_{\mathrm{bad}},p\rangle$ and a hypothetical proof $h$ of the universal sentence, apply $h$ at that host and class to obtain a cycle $z$, and feed $z$ to $p$ to obtain $\mathrm{False}$. This direction requires no classical axiom.

\paragraph{\textbf{Classical map (\texttt{negation\_to\_counterexample}).}}
From a function sending every purported proof of the universal sentence to $\mathrm{False}$, one does not constructively compute $D_{\mathrm{bad}}$. Extracting the existential witness uses \texttt{Classical.not\_forall}. The converse therefore depends on \texttt{Classical.choice}.

The kernel-checked lemmas in \texttt{Disproof/Type.lean} are \texttt{counterexample\_to\_negation}, \texttt{negation\_to\_counterexample}, and \texttt{equivalence}. Running \texttt{\#print axioms equivalence} reports \texttt{[propext, Classical.choice, Quot.sound]}. That theorem relates two types without constructing an inhabitant of either.

\section{The Miss Locus}

Write $L(D)$ for the rational Hodge conjecture in codimension two on a polarized fourfold $D$, $\Delta_{\mathrm{Hdg}}(D)$ for the primitive rational $(2,2)$ classes, and
\[
  \Delta_{\mathrm{miss}}(D)=\Delta_{\mathrm{Hdg}}(D)\setminus\operatorname{im}(\mathrm{cl}_D).
\]
By definition,
\[
  L(D)\;\Longleftrightarrow\;\Delta_{\mathrm{miss}}(D)=\emptyset.
\]
On a very general hypersurface fourfold of degree $d\ge 6$ in $\mathbb{P}^5$, the Lefschetz hyperplane theorem together with monodromy gives Noether--Lefschetz vanishing
\[
  H^4(X,\mathbb{Q})\cap H^{2,2}(X)=\mathbb{Q}\langle\omega^2\rangle,
\]
hence $\Delta_{\mathrm{Hdg}}(D)=0$ and $\Delta_{\mathrm{miss}}(D)=\emptyset$. The rational Hodge conjecture holds there because no primitive rational class remains to algebraize. A large complex Hodge number $h^{2,2}$ counts complex harmonic forms of type $(2,2)$, not rational Hodge classes.

A term of \texttt{RationalCounterexample} would require a \emph{special} fourfold with $\Delta_{\mathrm{Hdg}}\neq 0$ and $\Delta_{\mathrm{miss}}\neq\emptyset$. That host is not constructed in this note.

\section{Scope: What Is Asserted and What Is Not Asserted}

\begin{enumerate}
  \item The integral Hodge conjecture is false (classical literature counterexamples cited below).
  \item The two rational-negation types \texttt{RationalHodgeNegation} and \texttt{RationalCounterexample} are classically equivalent in Lean~4 with zero \texttt{sorry}.
  \item The constructive direction \texttt{counterexample\_to\_negation} uses no classical axiom.
  \item Neither rational-negation type is inhabited in this repository, nor is the universal rational sentence itself.
  \item This note does not claim the Clay Millennium Prize, and does not claim that the rational Hodge conjecture is either true or false.
\end{enumerate}

\begin{thebibliography}{9}
\bibitem{AH61}
M.~F. Atiyah and F.~Hirzebruch, \emph{Analytic cycles on complex manifolds}, Topology \textbf{1} (1961), 25--45.

\bibitem{Kol90}
J.~Koll\'ar, \emph{Trento examples}, in \emph{Classification of irregular varieties} (Trento, 1990), Lecture Notes in Math., vol.~1515, Springer, 1992, pp.~134--139.

\bibitem{Voi02}
C.~Voisin, \emph{Hodge Theory and Complex Algebraic Geometry I, II}, Cambridge Studies in Advanced Mathematics, Cambridge University Press, 2002.

\bibitem{Del06}
P.~Deligne, \emph{The Hodge Conjecture}, Clay Mathematics Institute Millennium Prize Problem description, 2000/2006.
\end{thebibliography}

\end{document}
